\section{Regime 1: Explicit Correction Single Stepping}\label{sec:regime1}

Recall the mismatch identity \eqref{eq:mismatchid}: forcing the fast equation (optimal step $\tau_1=\dto^{(1)}=\dx^2/(6D_1)$) onto a larger step $h$ re-involves the $u_{xxxx}$ term, and writing $h = \tau_1 + \varepsilon$ reduces its Group-1 coefficient to exactly $\frac{D_1^2\varepsilon}{2}$ per unit time (the per-step error carries the extra factor $h$). With this, we can formulate a first method of attempting to correspond any two coupled equations. 

For any two equations with distinct $\dto^{(1)} < \dto^{(2)}$, we evolve both equations by the shared step $h = \dto^{(2)}$ such that $u^{(2)}$ is optimally solved by nature, and the discretization of $u^{(1)}$ is evolved with the forced step $h = \tau_1 + \varepsilon$ such that the explicit mismatch correction term recovers order 4 accuracy. Namely,

\begin{lemma}\label{thm:mismatch}
  Forcing the time-step $h$ on the equation for species $i$, whose optimal step is $\tau_i:=\dto^{(i)}=\dx^2/(6D_i)$, write $h = \tau_i + \varepsilon$. The \textbf{explicit mismatch correction term} is

  \[
    C^{(i)}_{mismatch} = \frac{h D_i^2 \varepsilon}2 \delta_x^4 u^{(i), n}
  \]

  where $\delta_x^4=(\delta_x^2)^2$ is evaluated via a five-point stencil and $u^{(i), n}$ is the $i$-th species evaluated at time $n$.
\end{lemma}

Then,

\begin{theorem}[OTS-NIDC accuracy fixed indices with time-step forcing]\label{thm:cpl-r1-p1}
  For equation $u^{(i)}$ with forced time-step $h$, the scheme

  \begin{equation}\label{eq:coupledscheme1}
      u^{(i), n+1}_{j} = u^{(i), n}_{j} + h \left[ D_i \ddx u^{(i), n}_{j} + \sum_{\ell} R_{i\ell} u^{(\ell), n}_{j} \right] + C^{(1), (i)} + C^{(i)}_{mismatch}  
  \end{equation}

  has similar step error and global order as given by Theorem \ref{thm:p1}.
\end{theorem}


\subsection{Stability}

As we've established that the regime for individual equations remains the same, we prove stability by showing the correction factor described in Lemma \ref{thm:mismatch} is similarly a perturbation on the simple FE scheme. We follow the scalar analysis of \S\ref{sec:method}, now allowing the actual step $h$ to differ from the optimal step $\tau=\dto=\frac{\dx^2}{6D}$ (here $D=D_1$, the fast species). Write $h=\tau+\varepsilon$ and $r=\frac{h}{\tau}$; this is the $r$ of \S\ref{sec:method}, with $\mu=r/6$ and $a=\alpha h=r\eta/6$. We first establish stability of the mismatch-corrected diffusion step, then restore the reaction terms.

\begin{lemma}[Mismatch correction and amplification factor]
For $u_t=Du_{xx}$, the corrected update is
$u^{n+1}=u^n+hD\delta_x^2u^n
+\frac{hD^2\varepsilon}{2}\delta_x^4u^n$,
where $\delta_x^4=(\delta_x^2)^2$.
With $s=\sin^2\left(\frac{\kappa\dx}{2}\right)\in[0,1]$, the diffusion-only amplification factor of this update is
$G_{\mathrm{mis}}(s)
=1-\frac{2r}{3}s+\frac{2r(r-1)}{9}s^2$.
\end{lemma}

\begin{proof}
  By Lemma \ref{thm:mismatch}, cancellation requires adding $\frac{hD^2\varepsilon}{2}\delta_x^4u^n$.

The symbols of $\delta_x^2$ and $\delta_x^4$ are $-\frac{4s}{\dx^2}$ and $\frac{16s^2}{\dx^4}$. Since $\frac{hD}{\dx^2}=\frac{r}{6}$ and $\frac{hD^2\varepsilon}{\dx^4}=\frac{r(r-1)}{36}$ (substituting $h=r\tau$, $\varepsilon=(r-1)\tau$, $\tau=\dx^2/6D$), the update factor is
$1+hD\cdot\bigl(-\frac{4s}{\dx^2}\bigr)+\frac{hD^2\varepsilon}{2}\cdot\frac{16s^2}{\dx^4}
=G_{\mathrm{mis}}(s)
=1-\frac{2r}{3}s+\frac{2r(r-1)}{9}s^2$.
At $r=1$, the mismatch vanishes and this reduces to the diffusion part of the optimal-step factor \eqref{eq:factor} from \S\ref{sec:method}.
\end{proof}

Unlike the reaction corrections in \S\ref{sec:method}, the mismatch contribution need not vanish under refinement at fixed grid frequency: its symbol is $\frac{2r(r-1)}{9}s^2$. We therefore bound the complete factor rather than assume that the correction is uniformly small.

\begin{lemma}[Diffusive stability interval]
On an infinite or periodic uniform grid, the mismatch-corrected diffusion step satisfies $|G_{\mathrm{mis}}(s)|\le1$ for every $s\in[0,1]$ whenever $0<r\le4$.
\end{lemma}

\begin{proof}
The upper bound follows from
$G_{\mathrm{mis}}(s)-1
=\frac{2rs}{9}\left((r-1)s-3\right)\le0$:
for $r\le1$ the second factor is nonpositive; for $1<r\le4$ it is bounded by $(r-1)s\le r-1\le3$.

For the lower bound, completing the square gives
$G_{\mathrm{mis}}(s)
=\frac{2}{9}\left(rs-\frac{s+3}{2}\right)^2
+1-\frac{(s+3)^2}{18}$,
and the residual term satisfies $1-\frac{(s+3)^2}{18}\ge\frac19$ because $(s+3)^2\le16$ for $0\le s\le1$. Hence $\frac{1}{9}\le G_{\mathrm{mis}}(s)\le1$; the constant $\frac19$ is sharp, attained at $(r,s)=(2,1)$.

For $r>4$, the grid-scale factor is
$G_{\mathrm{mis}}(1)=1+\frac{2r(r-4)}{9}>1$.
Thus $r=4$ is the sharp boundary for this Fourier criterion.
\end{proof}

In particular, the correction preserves diffusive stability throughout the mismatch-corrected range $0<r\le4$ --- extending the uncorrected forward Euler ceiling $r\le3$ --- and in particular for shortened remainder steps $0<r<1$. We retain $r<4$ as the operating restriction (the strict form needed below) and use it to recover the reaction-growth bound from \S\ref{sec:method}.

For $u_t=Du_{xx}+\alpha u$, the reaction correction must also be evaluated at the actual step. The same Taylor expansion used in \S\ref{sec:method} gives
$C_{\mathrm{react}}u^n
=h^2\alpha D\delta_x^2u^n
+\frac{h^2\alpha^2}{2}u^n$.
At $h=\tau$, this is exactly the correction previously derived. Together with the mismatch term, it cancels the leading residual at a general step $h$.

\begin{remark}[Relation to the factors of \S\ref{sec:method}]\label{rem:gfac}
$G_{\mathrm{mis}}$ is the \emph{diffusion-only} building block and is not directly comparable to $G_{\mathrm{FE}}$, which includes the reaction. The full corrected factor derived in Theorem \ref{thm:r1-stab} below is named $G_{\mathrm{NIDC},h}$: it is the forced-step analogue of $G_{\mathrm{NIDC}}$, and plain forward Euler at the same step has $G_{\mathrm{FE}}(s)=\bigl(1-\frac{2r}{3}s\bigr)+a$, so
\[
G_{\mathrm{NIDC},h}(s)-G_{\mathrm{FE}}(s)
\;=\; \underbrace{\frac{2r(r-1)}{9}s^2}_{\text{mismatch}}
\;\underbrace{-\frac{2ra}{3}s\;+\;\frac{a^2}{2}}_{\text{reaction corrections at step }h},
\qquad a=\frac{r\eta}{6},
\]
which reduces at $r=1$ (where $a=\eta/6$) to exactly the perturbation $-\frac{\eta}{9}s+\frac{\eta^2}{72}$ of \S\ref{sec:method}. Off the optimum the two conventions differ: \S\ref{sec:method} evaluates the coefficients of $C^{(1)}$ (derived at $r=1$) at a general step, while the present section rescales the reaction corrections to $h$ and adds the mismatch term; the two coincide at $r=1$ and differ at $O(r-1)$ elsewhere (see the notation dictionary, Appendix~\ref{app:notation}).
\end{remark}

\begin{theorem}[Scalar stability with mismatch correction]\label{thm:r1-stab}
Let $\alpha\ge0$, and consider
$u^{n+1}
=u^n+h(D\delta_x^2u^n+\alpha u^n)
+\frac{hD^2(h-\tau)}{2}\delta_x^4u^n
+h^2\alpha D\delta_x^2u^n
+\frac{h^2\alpha^2}{2}u^n$.
On an infinite or periodic uniform grid, if $0<h<4\tau$, then its amplification factor satisfies $|G_{\mathrm{NIDC},h}(s)|\le e^{\alpha h}$ uniformly in $s\in[0,1]$.

Consequently, any sequence of admitted steps, including optimal substeps and shortened remainder steps, satisfies
$\|u^{n_f}\|_{\ell^2}\le e^{\alpha T}\|u^0\|_{\ell^2}$,
where $T=\sum_{n=0}^{n_f-1}h_n$ and $n_f$ is the final step index.
\end{theorem}

\begin{proof}
Set $a=\alpha h$ and $r=\frac{h}{\tau}$. Substituting the Fourier symbols gives
$G_{\mathrm{NIDC},h}(s)
=G_{\mathrm{mis}}(s)
+a\left(1-\frac{2r}{3}s\right)
+\frac{a^2}{2}$.
The preceding lemma gives $\frac19\le G_{\mathrm{mis}}(s)\le1$. Moreover, for $0<r<4$,
$-\frac{5}{3}<1-\frac{2r}{3}s\le1$.
Therefore, for the upper side,
$G_{\mathrm{NIDC},h}(s)\le1+a+\frac{a^2}{2}\le e^a$,
and for the lower side,
$G_{\mathrm{NIDC},h}(s)\ge\frac19-\frac{5a}{3}+\frac{a^2}{2}>-(1+a)\ge-e^a$,
where the middle inequality is $\frac{10}{9}-\frac{2a}{3}+\frac{a^2}{2}>0$, which holds for every $a\ge0$. Hence $|G_{\mathrm{NIDC},h}(s)|\le e^a$.

Over successive steps, the amplification factors multiply. Hence their product has magnitude at most
$\prod_{n=0}^{n_f-1}e^{\alpha h_n}=e^{\alpha T}$.
The discrete $\ell^2$ estimate follows from Parseval's identity.
\end{proof}

Thus the mismatch correction preserves the scalar growth allowance used in \S\ref{sec:method} throughout the admitted range $0<r<4$. In a subcycle, full steps have $r=1$ and any shortened remainder has $0<r<1$, so both satisfy the same estimate. This establishes the scalar building block; stability of off-diagonal reaction coupling and its interpolation still requires analysis of the complete coupled update.

\subsection{Advantage}
